%==============================================================================
% Resumen — APA 7: máximo 250 palabras + palabras clave
%==============================================================================
\newpage
\section*{Resumen}

La digitalización de los procesos organizacionales incrementa la dependencia de los sistemas informáticos y la criticidad de la gestión de incidentes. En muchas organizaciones medianas el registro inicial aún es manual, lo que genera demoras, errores de clasificación y uso ineficiente de recursos técnicos.

Esta investigación diseña, implementa y valida un sistema automatizado de mesa de ayuda que integra orquestación con N8N, un módulo de lenguaje natural en Python con FastAPI, inferencia con Gemini~3.6 Flash, persistencia PostgreSQL y telefonía con Twilio. Recibe solicitudes por correo electrónico, formulario web y teléfono, y deriva cada incidente a uno de los cinco sectores canónicos mediante un clasificador híbrido de reglas determinísticas y un modelo de lenguaje grande.

El estudio adopta un enfoque cuantitativo cuasiexperimental de muestras pareadas sobre los mismos 200 casos, con muestreo estratificado. El doble etiquetado independiente alcanzó un Kappa de Cohen de 0,87.

Los resultados evidencian una reducción estadísticamente significativa del tiempo medio de registro, de 61,6~s a 17,76~s (71,2~\%). La prueba de Wilcoxon arrojó \(W = 363\) y un tamaño del efecto \(r = 0{,}9639\). La exactitud estricta alcanzó 73,5~\% y el macro-F1 de cinco clases, 0,5207. El estado de la hipótesis principal se declara por sub-hipótesis: se confirma en la reducción del tiempo de registro, se refuta en el umbral de clasificación y la sub-hipótesis de intervención humana por caso se declara no contrastada por ausencia de registro primario.

Los hallazgos confirman la viabilidad técnica y económica; la operativa es condicional, porque el balance neto frente a las derivaciones erróneas resulta indeterminado. Se aporta una arquitectura reproducible y evidencia empírica en español rioplatense.

\bigskip
\noindent\textbf{Palabras clave:} automatización, mesa de ayuda, procesamiento de lenguaje natural, modelo de lenguaje grande, N8N, FastAPI, clasificación de tickets, español rioplatense, human in the loop

\newpage
\section*{Abstract}

The digitalization of organizational processes increases dependence on information systems and, with it, the criticality of incident management. In many medium-sized organizations, initial incident registration is still performed manually, introducing delays, classification errors, and inefficient use of technical resources.

This thesis designs, implements, and validates an automated help desk system integrating N8N workflow orchestration, a Python natural language processing module built with FastAPI, inference through the Gemini~3.6 Flash model, PostgreSQL persistence, and a telephony channel managed by Twilio. It receives requests through email, web form, and phone calls with automatic transcription, and routes each incident to one of the five canonical sectors using a hybrid classifier that combines deterministic rules with a large language model.

The study adopts a quantitative approach with a quasi-experimental paired-samples design over the same 200 cases, drawn by proportional stratified sampling. Independent double labeling reached a Cohen's Kappa of 0.87.

Results show a statistically significant reduction in mean registration time, from 61.6~s to 17.76~s (71.2~\%). The Wilcoxon signed-rank test yielded \(W = 363\) with an effect size \(r = 0.9639\). Strict classification accuracy reached 73.5~\% and the five-class macro-F1 was 0.5207. The state of the main hypothesis is declared per sub-hypothesis: it is confirmed in registration time, refuted in the classification threshold, and not tested due to the absence of a primary record for per-case human intervention.

The findings confirm technical and economic feasibility; operational feasibility is conditional, since the net balance against misrouted incidents remains undetermined. A reproducible architecture and empirical evidence in Rioplatense Spanish are provided.

\bigskip
\noindent\textbf{Keywords:} automation, help desk, natural language processing, large language model, N8N, FastAPI, ticket classification, Rioplatense Spanish, human in the loop

\newpage